% STATUS: VERIFIED / AUDITED
% AUDIT FILE: audits/A07_five_node_specialization.tex
% AUDIT VERDICT: VERIFIED
% LAYER VERDICTS: Fourier/trigonometry VERIFIED; ordering/selection VERIFIED;
%   frequency/coincidence VERIFIED; stability/critical-mode VERIFIED;
%   scope/dependency VERIFIED; overall VERIFIED.
% REQUIRED CORRECTIONS: none; no mathematical or wording revision.
% AUDITED VERSION: P07 v1.0; 2026-10-03; 24180 bytes; 555 lines.
% AUDITED FULL-FILE SHA-256:
%   fc10d257e8126b79a56cc7dd4c55b14b0266e250cbe40dbecb6c6a18658d09de
% AUDITED MATHEMATICAL-BODY SHA-256:
%   66a5b8f078cb0fc30f4d7c72c22a26b965fd8ded491f53450d604fe7a21ccdea
% CLOSURE DATE: 2026-10-03
% CURRENT VERSION: P07 v1.1; metadata-only closure.
% CURRENT FULL-FILE HASH: recorded separately in project_manifest.md and
%   audit_closure_report.txt to avoid a self-referential full-file hash.
% CURRENT MATHEMATICAL-BODY SHA-256:
%   66a5b8f078cb0fc30f4d7c72c22a26b965fd8ded491f53450d604fe7a21ccdea
% HASH SCOPE: bytes from the first line-anchored section command through EOF.
% BODY PRESERVATION: every mathematical-body byte is unchanged.
% STATUS NOTE: historical DRAFT/audit-pending body wording is preserved;
%   this header and the current manifest govern lifecycle status.
% VALIDATION PROVENANCE: A07 authenticated and exactly reproduced the
%   original 161 exact checks and 108 numerical cases and result files.
%   No historical validation suite is rerun for this closure.
% DOWNSTREAM: P09 cleared to begin, but not begun in this closure.
% DEPENDS ON: P03-P05 VERIFIED / AUDITED
% MODEL: canonical DM, n=5, q=2
% ROLE: five-node motivating specialization / exact corollary package
% NOVELTY ROLE: supporting specialization, not standalone golden-ratio novelty
% NUMERICAL DEPENDENCIES: none
% NEXT DEPENDENCIES: P11-P14 when five-node specialization is used
% VERSION: P07 v1.1; 2026-10-03; metadata-only A07 audit closure
% DIRECT DEPENDENCY: P05; P03/P04 and P01/P02 are inherited prerequisites.
% CONTRACT: 00_canonical_model.tex; 01_notation_and_assumptions.tex;
%           shared_macros.tex. No P08 theorem dependency is introduced.
% BASELINE: cyclic_network_project_P08_audited.zip; 118 authenticated members.
% CLAIMS: C051, C052, C155 (interpretation rejected); P07-D1--P07-D6.
% SCOPE: exact spectral algebra at a specified positive synchronous equilibrium;
% no parameter-path first instability, Hopf theorem, or nonlinear conclusion.

\section{P07: Five-node golden-ratio specialization}
\label{sec:P07}
\noindent\textbf{Status: DRAFT; independent A07 audit pending.}
This is a corollary package for the audited P03--P05 demand-modulated
results. The direct dependency remains $\mathrm{P05}\to\mathrm{P07}$;
P03--P04 and the P01--P02 equilibrium foundation are inherited.
P08 is present in the current project but supplies no theorem used here.

Fix $a,k,c,g,h>0$ in the canonical model
\begin{equation}
 \dot x_i=ax_i-kx_i^2+\frac{cx_{i-1}}{1+x_i}
       -\frac{gx_i}{1+x_{i+2}}-hx_ix_{i-2},\qquad i\in\Z_5.
 \label{eq:P07-model}
\end{equation}
Let $x^*\one$, $x^*>0$, be a specified synchronous equilibrium. Write
\begin{equation}
 s=1+x^*>1,\quad C=\frac cs>0,\quad G=\frac{gx^*}{s^2}>0,
 \quad H=hx^*>0,\quad T=G-H,\quad S=G+H.
 \label{eq:P07-coefficients}
\end{equation}
Here $s$ denotes the denominator, not P02's scalar state variable;
$S$ is local to P07 and denotes effective DM coefficients, not the
physical sum $g+h$ used locally in P08. We use $D_5:=H-G=-T$,
leaving P02's discriminant $D$ unchanged. The symbols $\omega_1,\omega_2$
below are signed imaginary parts, not a reserved positive Hopf frequency.
By P02 and P03,
\begin{align}
 a-(k+h)x^*+\frac{c-g}{s}&=0,\label{eq:P07-equilibrium}\\
 d&=a-2kx^*-\frac{cx^*}{s^2}-\frac gs-hx^*
   =-kx^*-C-\frac{Cx^*}{s}.\label{eq:P07-d}
\end{align}
The last equality uses equilibrium and will only be used when stated.

\subsection{Exact fifth-root constants and Fourier spectrum}
\label{slot:P07-1}
\begin{lemma}[Golden-ratio trigonometry]
\label{lem:P07-trig}
Set
\begin{equation}
 \phi=\frac{1+\sqrt5}{2},\qquad
 u_1=\cos\frac{2\pi}{5},\qquad u_2=\cos\frac{4\pi}{5},\qquad
 \sigma=\sin\frac\pi5.
 \label{eq:P07-constants}
\end{equation}
Then
\begin{gather}
 \phi^2=\phi+1,\quad \phi^{-1}=\phi-1,\quad
 \phi^{-2}=2-\phi=\frac{3-\sqrt5}{2},\quad
 \phi+\phi^{-1}=\sqrt5,\label{eq:P07-phi}\\
 u_1=\frac{\sqrt5-1}{4}=\frac1{2\phi},\qquad
 u_2=-\frac{\sqrt5+1}{4}=-\frac\phi2,\qquad
 u_1-u_2=\frac{\sqrt5}{2},\label{eq:P07-cosines}\\
 \sigma=\frac{\sqrt{10-2\sqrt5}}4>0,\qquad
 \sin\frac{2\pi}{5}=\phi\sigma
       =\frac{\sqrt{10+2\sqrt5}}4.\label{eq:P07-sines}
\end{gather}
\end{lemma}
\begin{proof}
The quadratic identity for $\phi$ follows by squaring its definition;
division by $\phi>0$ and another multiplication give the remaining
identities in \eqref{eq:P07-phi}. Let $z=e^{2\pi\ii/5}$.
Since $1+z+z^2+z^3+z^4=0$, real parts give
$1+2u_1+2u_2=0$. Also $u_2=2u_1^2-1$, so
$4u_1^2+2u_1-1=0$. As $0<2\pi/5<\pi/2$, the positive root
is $u_1=(\sqrt5-1)/4$; substitution gives $u_2$ and their difference.
Moreover $\cos(\pi/5)=-\cos(4\pi/5)=\phi/2$.
Thus $\sigma^2=1-\phi^2/4=(5-\sqrt5)/8$, with the positive
square root required by $0<\pi/5<\pi/2$. The double-angle formula
then gives $\sin(2\pi/5)=2\sigma(\phi/2)=\phi\sigma$;
squaring and using $\phi^2=\phi+1$ gives the second radical.
\end{proof}

With $\theta_m=2\pi m/5$, all phase signs needed below are displayed
explicitly. Each row follows from periodicity and the quadrant signs:
\begin{center}
\begin{tabular}{@{}c r r r r@{}}
\toprule
$m$&$\cos\theta_m$&$\sin\theta_m$&$\cos(2\theta_m)$&$\sin(2\theta_m)$\\
\midrule
0&$1$&$0$&$1$&$0$\\
1&$u_1$&$\phi\sigma$&$u_2$&$\sigma$\\
2&$u_2$&$\sigma$&$u_1$&$-\phi\sigma$\\
3&$u_2$&$-\sigma$&$u_1$&$\phi\sigma$\\
4&$u_1$&$-\phi\sigma$&$u_2$&$-\sigma$\\
\bottomrule
\end{tabular}
\end{center}
There are exactly three representative classes: $j=0$ with labels
$\{0\}$, $j=1$ with labels $\{1,4\}$, and $j=2$ with labels
$\{2,3\}$. Odd dimension excludes a Nyquist singleton. The two labels
of a conjugate pair are one real-part class, even when their eigenvalues
are distinct.

\begin{theorem}[Exact five-node Fourier spectrum]
\label{thm:P07-spectrum}
The five eigenvalues, with their Fourier labels and algebraic
multiplicities retained, are
\begin{align}
 \lambda_0&=d+C+T,\label{eq:P07-spectrum}\\
 \lambda_1&=d+\frac{C}{2\phi}-\frac{\phi T}{2}
                +\ii\sigma(S-\phi C),\nonumber\\
 \lambda_2&=d-\frac{\phi C}{2}+\frac{T}{2\phi}
                -\ii\sigma(C+\phi S),\nonumber\\
 \lambda_4&=\overline{\lambda_1},\qquad
 \lambda_3=\overline{\lambda_2}.\nonumber
\end{align}
In particular the three real-part classes are
\begin{equation}
 r_0=d+C+T,\qquad
 r_1=d+\frac{C}{2\phi}-\frac{\phi T}{2},\qquad
 r_2=d-\frac{\phi C}{2}+\frac{T}{2\phi}.
 \label{eq:P07-real}
\end{equation}
Here $r_1=\Ree\lambda_1=\Ree\lambda_4$ and
$r_2=\Ree\lambda_2=\Ree\lambda_3$.
\end{theorem}
\begin{proof}
Start with the audited P04 spectrum, Theorem~\ref{thm:P04-spectrum}:
\[
 \lambda_m=d+Ce^{-\ii\theta_m}+Ge^{2\ii\theta_m}
                         -He^{-2\ii\theta_m}.
\]
Expansion into sine and cosine, term by term, yields
\[
 \lambda_m=d+C\cos\theta_m+T\cos(2\theta_m)
            +\ii[-C\sin\theta_m+S\sin(2\theta_m)].
\]
In particular $-He^{-2\ii\theta_m}$ contributes
$-H\cos(2\theta_m)+\ii H\sin(2\theta_m)$: the real coefficient
is $G-H$, whereas the imaginary coefficient is $G+H$.
Insert rows 0, 1 and 2 of the phase table and
\eqref{eq:P07-cosines}--\eqref{eq:P07-sines}. Rows 4 and 3 have
unchanged cosines and opposite sines, proving the stated conjugacy.
As a separate real-part check, P05's quadratic
$R(u)=d-T+Cu+2Tu^2$ gives $R(1)=d+C+T$.
The identities $2u_1^2-1=u_2$ and $2u_2^2-1=u_1$ give
$R(u_1)=d+Cu_1+Tu_2$ and $R(u_2)=d+Cu_2+Tu_1$,
exactly \eqref{eq:P07-real}.
\end{proof}

\subsection{Complete ordering and the dominant-mode switch}
\label{slot:P07-2}
\begin{theorem}[Complete ordering of the three real-part classes]
\label{thm:P07-ordering}
Set $t=T/C$, which is well defined since $C>0$. The exact differences are
\begin{align}
 r_1-r_2&=\frac{\sqrt5}{2}(C-T)
              =\frac{\sqrt5 C}{2}(1-t),\label{eq:P07-gaps}\\
 r_0-r_1&=\frac{\sqrt5\phi}{2}
                    \left(T+\frac{C}{\phi^2}\right)
              =\frac{\sqrt5\phi C}{2}(t+\phi^{-2}),\nonumber\\
 r_0-r_2&=\frac{\sqrt5}{2\phi}(T+\phi^2 C)
              =\frac{\sqrt5 C}{2\phi}(t+\phi^2).\nonumber
\end{align}
Consequently the following table is exhaustive, including all equalities:
\begin{center}
\renewcommand{\arraystretch}{1.15}
\begin{tabular}{@{}l l l@{}}
\toprule
Condition&Ordering&Tie type\\
\midrule
$t>1$&$r_0>r_2>r_1$&none\\
$t=1$&$r_0>r_1=r_2$&subdominant\\
$-\phi^{-2}<t<1$&$r_0>r_1>r_2$&none\\
$t=-\phi^{-2}$&$r_0=r_1>r_2$&dominant\\
$-\phi^2<t<-\phi^{-2}$&$r_1>r_0>r_2$&none\\
$t=-\phi^2$&$r_1>r_0=r_2$&subdominant\\
$t<-\phi^2$&$r_1>r_2>r_0$&none\\
\bottomrule
\end{tabular}
\end{center}
Ordering is independent of $d$. No triple tie of real-part classes occurs.
\end{theorem}
\begin{proof}
Subtract \eqref{eq:P07-real}. The first gap is
$(\phi+\phi^{-1})(C-T)/2$. The second is
$C(1-1/(2\phi))+T(1+\phi/2)$; the third exchanges these
coefficients. Using \eqref{eq:P07-phi},
\[
 1+\frac\phi2=\frac{\sqrt5\phi}{2},\qquad
 1-\frac1{2\phi}=\frac{\sqrt5}{2\phi},\qquad
 \frac{1-1/(2\phi)}{1+\phi/2}=\phi^{-2}.
\]
This proves every factor in \eqref{eq:P07-gaps}. All prefactors are
positive, and $-\phi^2<-\phi^{-2}<1$. Starting above $1$, the
signs of $(r_1-r_2,r_0-r_1,r_0-r_2)$ are $(-,+,+)$.
They are $(0,+,+)$ at $1$, $(+,+,+)$ down to $-\phi^{-2}$,
$(+,0,+)$ there, $(+,-,+)$ down to $-\phi^2$,
$(+,-,0)$ there, and $(+,-,-)$ below it. These signs give
exactly the seven rows. Distinct gap zeros exclude a triple class tie.
The common summand $d$ cancels from each difference; it is not an
independent physical parameter by this observation.
\end{proof}

\begin{corollary}[Exact golden-ratio selection threshold]
\label{cor:P07-selection}
The dominant representative set and the full selected Fourier-label set are
\begin{equation}
\begin{array}{c|c|c}
\text{Condition}&\text{Representatives}&\text{Labels}\\\hline
D_5/C<\phi^{-2}&\{0\}&\{0\}\\
D_5/C=\phi^{-2}&\{0,1\}&\{0,1,4\}\\
D_5/C>\phi^{-2}&\{1\}&\{1,4\}.
\end{array}
\label{eq:P07-selection}
\end{equation}
Pair $\{2,3\}$ is never dominant at a canonical positive synchronous
 equilibrium. The unique dominant-class switch is
\begin{equation}
 \boxed{H-G=\frac{C}{\phi^2},\qquad
         \frac{H-G}{C}=\frac{3-\sqrt5}{2}.}
 \label{eq:P07-switch}
\end{equation}
\end{corollary}
\begin{proof}
The sign of $r_0-r_1$ gives the three cases. For a direct exclusion
of class 2, if $t<1$ then $r_1>r_2$ by the first gap; if $t\ge1$
then $r_0>r_2$ by the third gap. Thus class 2 cannot maximize even
at either subdominant equality. This argument compares all three
classes and does not rely solely on P05's reachability list.
\end{proof}

\paragraph{Check against P05's exact cells.}
In dimension five,
\begin{equation}
 \beta_0=\frac{1+u_1}{2}=\frac{\phi^2}{4},\qquad
 \beta_1=\frac{u_1+u_2}{2}=-\frac14,\qquad
 \alpha_0=\beta_0^{-1}=\frac4{\phi^2}=6-2\sqrt5.
 \label{eq:P07-midpoints}
\end{equation}
Indeed $2+\phi^{-1}=\phi+1=\phi^2$ and
$\phi^{-1}-\phi=-1$. In the concave regime $D_5>0$, P05's
vertex is $u_*=C/(4D_5)>0$. It passes $\beta_0$ at
\eqref{eq:P07-switch} and never reaches $\beta_1$.
Theorem~\ref{thm:P05-selection} therefore agrees exactly with the
three-class calculation. In P05's variables,
\begin{equation}
 \chi=\frac{4(H-G)}C,\qquad
 \begin{cases}
 \chi<4/\phi^2 &: \{0\}\text{ selected},\\
 \chi=4/\phi^2 &: \{0,1\}\text{ selected},\\
 \chi>4/\phi^2 &: \{1\}\text{ selected}.
 \end{cases}
 \label{eq:P07-chi}
\end{equation}
The current P05 defines $\chi$ for all real signs, since $C>0$,
so \eqref{eq:P07-chi} is also valid when $H\le G$:
then $\chi\le0$ and only class 0 is selected. The concave vertex
and reciprocal-midpoint argument above are used only when $H>G$.

\paragraph{Canonical switch surface and the equilibrium constraint.}
Since $D_5=x^*(h-g/s^2)$, solving \eqref{eq:P07-switch} gives
\begin{equation}
 \boxed{h=\frac{g}{(1+x^*)^2}
            +\frac{c}{\phi^2(1+x^*)x^*}.}
 \label{eq:P07-canonical-switch}
\end{equation}
The divisions are valid because $x^*,s,\phi,C$ are positive.
The equilibrium relation \eqref{eq:P07-equilibrium} must hold
simultaneously. Thus this equation describes a switch on the positive
branch-state manifold; $x^*$ cannot be prescribed independently of
a fixed physical tuple. The corresponding strict inequalities in $h$
at that equilibrium give uniform or pair-1 selection, respectively.
P05 already supplies regularity and local realizability of this switch.

\paragraph{The other golden-ratio equality surfaces.}
At $t=1$, equivalently $G-H=C$, one has $r_1=r_2$ but
$r_0-r_1=r_0-r_2=5C/2>0$. At $t=-\phi^2$, equivalently
$H-G=\phi^2C$, one has $r_0=r_2$ but
$r_1-r_0=r_1-r_2=5\phi C/2>0$.
These values follow from \eqref{eq:P07-gaps} and $\phi^2=\phi+1$.
Both are subdominant real-part ties, not dominant-mode switches.
No equality of complex numerical eigenvalues is implied by either tie.

\subsection{Frequencies, intersection, and numerical eigenvalue equality}
\begin{proposition}[Signed imaginary parts and frequency degeneracy]
\label{prop:P07-frequency}
At positive synchrony,
\begin{equation}
 \omega_1:=\Imm\lambda_1=\sigma(S-\phi C),\qquad
 \omega_2:=\Imm\lambda_2=-\sigma(C+\phi S)<0.
 \label{eq:P07-frequency}
\end{equation}
Pair $\{2,3\}$ is always nonreal. Pair $\{1,4\}$ is nonreal
if and only if $S\ne\phi C$. At the spectral locus
\begin{equation}
 \boxed{G+H=\phi C,}\label{eq:P07-frequency-locus}
\end{equation}
one has $\lambda_1=\lambda_4=r_1$, a repeated real numerical
 eigenvalue represented by two Fourier labels. All numerical eigenvalues
are semisimple; any further coincident label is counted separately below.
\end{proposition}
\begin{proof}
The imaginary terms were derived directly in Theorem~\ref{thm:P07-spectrum}.
Since $\sigma,C,S,\phi>0$, $\omega_2$ is strictly negative, while
$\omega_1=0$ is equivalent to $S=\phi C$. Conjugacy then proves
all reality assertions. Semisimplicity is
Proposition~\ref{prop:P04-multiplicity}, which does not assume distinct
 numerical eigenvalues. These assertions classify spectra only.
\end{proof}

\begin{proposition}[Intersection with the dominant-mode switch]
\label{prop:P07-intersection}
The switch and the frequency-degeneracy locus intersect in the strictly
positive coefficient domain exactly at
\begin{equation}
 \boxed{\frac GC=\frac1\phi,\qquad \frac HC=1.}
 \label{eq:P07-intersection}
\end{equation}
They define a nonempty codimension-two spectral locus in $(C,G,H)>0$,
and locally on the canonical positive equilibrium branch. At that locus
\begin{equation}
 \lambda_0=\lambda_1=\lambda_4=d+\frac C\phi
   =-kx^*-\frac{Cx^*}{s}-\frac{C}{\phi^2}<0.
 \label{eq:P07-triple}
\end{equation}
This numerical eigenvalue is semisimple of multiplicity three, and
$r_2<r_0$. In particular the full linearization is Hurwitz there.
This is a spectral codimension statement, not a bifurcation classification.
\end{proposition}
\begin{proof}
Adding and subtracting $H-G=C\phi^{-2}$ and $G+H=\phi C$ gives
\[
 \frac GC=\frac{\phi-\phi^{-2}}2=\phi-1=\phi^{-1},\qquad
 \frac HC=\frac{\phi+\phi^{-2}}2=1.
\]
Both are positive. The gradients in $(G,H)$ of
$H-G-C/\phi^2$ and $G+H-\phi C$ have determinant $-2$, so the
two linear conditions are independent. For a canonical witness take
\begin{equation}
 x^*=1,\qquad (a,k,c,g,h)=(\sqrt5,1,2,4/\phi,1).
 \label{eq:P07-intersection-witness}
\end{equation}
The equilibrium bracket equals $\sqrt5-2+(2-4/\phi)/2=0$,
using $1+2/\phi=\sqrt5$. The effective coefficients are
$C=1,G=\phi^{-1},H=1$.

For the local branch assertion use P05's verified equilibrium chart
$(x,k,c,g,h)\mapsto((k+h)x-(c-g)/(1+x),k,c,g,h)$.
Its determinant is positive wherever $H-G>0$; the intersection lies
in that regime. At fixed $x,k,c$, the two constraint gradients in
$(g,h)$ have determinant $-2x^2/(1+x)^2\ne0$.
The chart therefore transfers their independent codimension-two
conditions locally to physical parameters on the specified branch.
There is no claim about a parameter-path crossing or an unfolding.

At the intersection, conjugacy, $\omega_1=0$, and $r_0=r_1$
give the triple value. Substitute \eqref{eq:P07-d} and
$1-\phi^{-1}=\phi^{-2}$ for its strict negativity.
The ordering table gives $r_2<r_0$; $\omega_2<0$ keeps labels 2 and 3
distinct from the real value and from each other. P04 counts precisely
three independent Fourier eigenvectors for the triple value.
\end{proof}

\begin{proposition}[Complete numerical coincidence classification]
\label{prop:P07-coincidences}
Under the stated positive-equilibrium assumptions, the only numerical
 eigenvalue coincidences are as follows:
\begin{enumerate}[label=(\roman*),leftmargin=*]
\item If $S\ne\phi C$, all five numerical eigenvalues are distinct.
\item If $S=\phi C$ and $T\ne-C/\phi^2$, only
$\lambda_1=\lambda_4$ coincide; their real value has multiplicity two.
\item If $S=\phi C$ and $T=-C/\phi^2$, exactly
$\lambda_0=\lambda_1=\lambda_4$ coincide as in
\eqref{eq:P07-triple}; their value has multiplicity three.
\end{enumerate}
All multiplicities are algebraic and geometric. Labels 2 and 3 always
give distinct, simple numerical eigenvalues, sharing no value with any
other label. Every nonreal pair-1 eigenvalue is simple.
\end{proposition}
\begin{proof}
Within pairs, \eqref{eq:P07-frequency} proves that labels 1 and 4
coincide exactly when $S=\phi C$, while labels 2 and 3 never coincide.
The real $\lambda_0$ cannot equal either nonreal pair-2 eigenvalue.
Equality between $\lambda_0$ and a pair-1 label requires both
$r_0=r_1$ and $\omega_1=0$, exactly the intersection just analyzed.

A cross-pair numerical equality requires $r_1=r_2$, hence $T=C$.
Then $G=H+C$ and $S=C+2H>C$. At this real-part tie,
\begin{align*}
 \omega_1-\omega_2
    &=\sigma[(1+\phi)S-(\phi-1)C]>0,\\
 -\omega_2-\omega_1
    &=\sigma[(1+\phi)C+(\phi-1)S]>0.
\end{align*}
For the first inequality use $S>C$ and
$(1+\phi)-(\phi-1)=2$; the second has strictly positive summands.
Thus $\omega_2<\omega_1<-\omega_2$, so
$\omega_1$ is neither $\omega_2$ nor $-\omega_2$; conjugating
excludes all four possible cross-pair equalities. This exhausts all
unordered label pairs. P04 then gives the asserted multiplicities.
\end{proof}
The frequency-degeneracy locus itself has $-\phi<t<\phi$, since
$G,H>0$ imply $|T|<S$. It can intersect the subdominant tie $t=1$
without a cross-pair coincidence, but cannot intersect $t=-\phi^2$.
Equal real parts, repeated real values, and distinct nonreal conjugates
are consequently kept separate throughout.

\subsection{Exact stability test and critical-mode ledger}
\begin{corollary}[Piecewise five-node linear stability]
\label{cor:P07-stability}
The linearization is exponentially asymptotically stable exactly when
\begin{equation}
 \max\{r_0,r_1,r_2\}<0
 \quad\Longleftrightarrow\quad
 \begin{cases}
 r_0<0,&t>-\phi^{-2},\\
 r_0=r_1<0,&t=-\phi^{-2},\\
 r_1<0,&t<-\phi^{-2}.
 \end{cases}
 \label{eq:P07-stability}
\end{equation}
The middle row means that the common value must be negative. At canonical
positive equilibria on that tie it is automatically negative:
\begin{equation}
 r_0=r_1=-kx^*-\frac{Cx^*}{s}-\frac{C}{\phi^2}<0.
 \label{eq:P07-stable-tie}
\end{equation}
\end{corollary}
\begin{proof}
P04 Corollary~\ref{cor:P04-stability} gives the strict all-mode test;
Theorem~\ref{thm:P07-ordering} gives its piecewise maximum.
At the dominant tie, substitute $T=-C/\phi^2$ and
\eqref{eq:P07-d} into $r_0=d+C+T$.
\end{proof}
If the maximum is zero, P04 implies that the normal linear system is
Lyapunov stable but not asymptotically stable. This does not settle
nonlinear stability or establish a bifurcation. A strictly positive
maximum means linear instability. The dominant tie cannot itself be a
zero-abscissa boundary in the canonical domain, by
\eqref{eq:P07-stable-tie}.

\begin{proposition}[Fixed-state critical-mode ledger]
\label{prop:P07-ledger}
Only the following classes can attain the spectral abscissa; this is a
fixed-state eligibility ledger, not a first-instability theorem.
\begin{center}
\begin{tabular}{@{}l l l@{}}
\toprule
Class&Eligibility for the maximum&Signed imaginary part\\
\midrule
$0$&$t\ge-\phi^{-2}$&$0$\\
$\{1,4\}$&$t\le-\phi^{-2}$&$\pm\sigma(S-\phi C)$\\
$\{2,3\}$&never&$\mp\sigma(C+\phi S)$\\
\bottomrule
\end{tabular}
\end{center}
At a zero-abscissa state, the inequalities in the first two rows must
be strict, because the shared boundary is strictly stable.
If $r_2=0$, some other class has strictly positive real part, so pair 2
cannot be the sole or tied controlling class of a stability boundary.
No oriented parameter path, prior stability, or crossing is asserted.
\end{proposition}
\begin{proof}
The eligibility and exclusions are precisely
Corollary~\ref{cor:P07-selection}. At every point at least one of
$r_0,r_1$ exceeds $r_2$ strictly. The zero-abscissa statements follow
from that strict comparison and \eqref{eq:P07-stable-tie}.
\end{proof}

\subsection{Realizability checks, downstream inputs, and scope}
\paragraph{All ordering regions have canonical representatives.}
For any $C,G,H>0$, the following single construction gives a canonical
positive equilibrium with those exact effective coefficients:
\begin{equation}
 x^*=1,\quad k=C,\quad c=2C,\quad g=4G,\quad h=H,
 \quad a=H+2G>0,\quad d=-\frac{5C}{2}.
 \label{eq:P07-realization}
\end{equation}
Indeed \eqref{eq:P07-equilibrium} becomes
$H+2G-(C+H)+(2C-4G)/2=0$. Every original component equals
$x^*$ times that zero bracket. For any real $t$, take $C=1$,
$G=1+\max(t,0)$ and $H=1+\max(-t,0)$. Then $G-H=t$ and both
coefficients are positive. Thus no row of the ordering table is merely
an unattached-coefficient possibility. This elementary five-node check
uses no MB comparison or new general selection theorem.

\paragraph{Exact inputs available to P11.}
The unique nonuniform dominant class, when present, is $\{1,4\}$.
Its signed imaginary part is $\sigma(S-\phi C)$, and its exact
nonzero-frequency condition is $S\ne\phi C$. Whenever pair 1 is
maximal, $r_2<r_1$ strictly, including the uniform/pair-1 dominant tie.
When pair 1 is uniquely dominant, $r_0<r_1$ as well.
Every nonreal pair-1 eigenvalue is simple; pair 2 never coincides
numerically with another label. The only uniform/pair-1 numerical
coincidence is the strictly stable triple value
\eqref{eq:P07-triple}; more generally the entire dominant tie is
strictly stable. These are exact spectral inputs, without a Hopf
existence or first-crossing conclusion. No transversality, crossing
derivative, Lyapunov coefficient, normal form, periodic orbit, or rotating
wave is analyzed here; P09--P18 remain unbegun by this module.

\begin{remark}[Golden ratio is network geometry]
\label{rem:P07-geometry}
The constants $\phi$ and $\sigma$ arise solely from the fifth roots of
unity and the five-node Fourier grid. They establish no physiological
golden-ratio law, biological optimum, medical significance, or empirical
TCM validation. P07 is a supporting specialization of the audited general
program, not standalone golden-ratio novelty. The frozen scope and
novelty documents need no amendment.
\end{remark}

\paragraph{Supplementary corroboration.}
The separate \path{P07_validation_support.zip} contains exact algebraic
checks, direct $5\times5$ matrix comparisons at positive canonical
 equilibria, current-input hashes, and build/preservation records.
Checks cover all seven ordering rows, both sides of the dominant switch,
the other two ties, and the frequency locus and its neighborhood.
The analytic proofs above do not depend on a finite check, floating-point
equality, compilation, or numerical continuation. P07 remains DRAFT.

\subsection{Priority checks for the independent P07 auditor}
\begin{enumerate}[leftmargin=*]
\item Reconstruct the fifth-root cosine quadratic, positive sine radicals,
and every use of $\phi^2=\phi+1$; check all five phase-table rows.
\item Start from P04's oriented shift phases and check the minus sign in
$-He^{-2\ii\theta_m}$. The real coefficient is $T=G-H$; the
imaginary coefficient is $S=G+H$. Verify both signed frequencies.
\item Recompute all three factorized gaps and reconstruct all seven
ordering rows. Distinguish the dominant threshold $-\phi^{-2}$ from
the subdominant threshold $-\phi^2$ and the other subdominant tie $t=1$.
\item Check all three classes directly, the impossibility of pair-2
dominance, the midpoint constants, the all-sign P05 convention for
$\chi$, and the canonical switch together with the equilibrium relation.
\item Check $S=\phi C$, the positive intersection $G/C=\phi^{-1},H/C=1$,
the codimension gradients and branch regularity, and the strictly
negative triple value. Count numerical multiplicity rather than
conjugate labels or tied real parts.
\item Exhaust all label-pair coincidences, particularly the cross-pair
real-part tie $T=C$, where $S>C$ is essential. Verify the independent
strict negativity of the whole dominant tie, not only its intersection.
\item Check the piecewise Hurwitz test and the critical-mode ledger;
retain the distinction between a fixed-state maximizer and a first
instability along a path. No Hopf or other bifurcation claim is licensed.
\item Authenticate the current P03--P05 and audit closures, preserve
$\mathrm{P05}\to\mathrm{P07}$ without a P08 edge, and keep P07 DRAFT
until a separate exact-version audit is completed. Computational
corroboration and a successful build are not that audit.
\end{enumerate}
